\subsection{Power Series and Approximation}

% **Role:** convergence of functions and analytic approximation.
%
% Power series, Taylor series, analytic functions, Weierstrass approximation, Stone--Weierstrass.


Power series are controlled infinite polynomials, and approximation is uniform closure. This subsection covers radii, Taylor structure, Abel behavior, Weierstrass, Stone--Weierstrass, and elementary analytic functions.

\begin{motivation} \label{mot:power-series-as-controlled-infinite-polynomials} \hlt{Power Series as Controlled Infinite Polynomials}\\
Power-series manipulations rest on uniform convergence: Definition \ref{def:series-and-partial-sums}, Theorems \ref{thm:power-series-radius-convergence}, \ref{thm:weierstrass-m-test}, and \ref{thm:uniform-derivative-convergence}.
\end{motivation}

\begin{definition} \label{def:power-approx-fourier-core-source-todo}\label{def:power-approx-core-objects} \hlt{Core Objects: Power Series, Approximation, and Gamma}
\begin{enumerate}[label=(\roman*)]
\setlength{\itemsep}{0pt}
\item A real \hlt{power series about $a$} is $\sum\limits_{n=0}^{\infty}c_n(x-a)^n$; a local representation of $f$ by one makes $f$ \hlt{analytic} at $a$.
\item A \hlt{polynomial approximant} to $f$ on a compact set $K$ is a polynomial $P$ chosen so that $\norm{f-P}_{\infty,K}$ is small.
\item A function algebra is a collection closed under addition, multiplication, and scalar multiplication; Stone--Weierstrass asks when such algebras are uniformly dense in $C(K)$.
\item For $x>0$, the \hlt{gamma function} is $\Gamma(x)\defeq\int_0^\infty t^{x-1}e^{-t}\,dt$ when the improper integral converges.
\end{enumerate}
\hlt{Key Insight}: Power series are local analytic objects; approximation theorems turn algebraic families into dense function families; $\Gamma$ extends factorials by integration.
\end{definition}

\begin{definition} \label{def:power-series-radius-basic} \hlt{Power Series and Radius of Convergence}\\
Given a complex sequence $(c_n)$, the series $\sum\limits_{n=0}^{\infty}c_nz^n$ is a \hlt{power series}. If $\alpha\defeq\limsup\limits_{n\to\infty}\sqrt[n]{\abs{c_n}}$, define $R\defeq1/\alpha$, with $R=+\infty$ when $\alpha=0$ and $R=0$ when $\alpha=+\infty$. The number $R$ is the \hlt{radius of convergence}.
\end{definition}

\begin{theorem} \label{thm:power-series-radius-convergence} \hlt{Convergence Inside the Circle}\\
For $\sum c_nz^n$, radius $R$: converge if $\abs{z}<R$, diverge if $\abs{z}>R$; boundary $\abs{z}=R$ is undecided.
\begin{proof}
For $a_n=c_nz^n$, $\limsup\limits_{n\to\infty}\sqrt[n]{\abs{a_n}}=\abs{z}\limsup\limits_{n\to\infty}\sqrt[n]{\abs{c_n}}=\frac{\abs{z}}{R}$, with extended meaning if $R=0$ or $R=+\infty$.\\
If $\abs{z}<R$, this value is $<1$; if $\abs{z}>R$, it is $>1$. Apply Theorem \ref{thm:root-test}.\\
\hlt{Structural Commentary}: Radius controls the interior and exterior, not the boundary.\\
\hlt{Key Insight}: The root test is the radius-of-convergence theorem in disguise.
\end{proof}
\end{theorem}

\begin{theorem} \label{thm:power-series-boundary-dirichlet} \hlt{A Boundary Test for Power Series}\\
If $\sum c_nz^n$ has radius $1$, $c_n\downarrow0$, and $c_n\geq0$, then it converges for $\abs{z}=1$, except possibly at $z=1$.
\begin{proof}
For $\abs{z}=1$ and $z\neq1$, put $a_n=z^n$ and $b_n=c_n$.\\
The partial sums of $\sum z^n$ are bounded because $\abs{\sum\limits_{m=0}^{n}z^m}=\abs{\frac{1-z^{n+1}}{1-z}}\leq\frac{2}{\abs{1-z}}$. Theorem \ref{thm:dirichlet-test-series} applies.\\
\hlt{Structural Commentary}: Boundary convergence away from $z=1$ comes from bounded geometric partial sums.\\
\hlt{Key Insight}: Monotone coefficients can make the unit-circle boundary converge without absolute convergence.
\end{proof}
\end{theorem}

\begin{perspective} \label{per:centered-power-series-translation} \hlt{Centered Series Are Translated Series}\\
$\sum c_n(x-a)^n$ is $\sum c_ny^n$ with $y=x-a$. Its radius is $R$ from Definition \ref{def:power-series-radius-basic}: absolute convergence for $\abs{x-a}<R$, divergence for $\abs{x-a}>R$, and separate endpoint tests for $\abs{x-a}=R$.
\end{perspective}

\begin{figure}[H]
  \centering
  \begin{tikzpicture}[
  >=Stealth,
  every node/.style={font=\scriptsize},
  axis/.style={->, gray!70, line width=0.5pt},
  conv/.style={blue!70!black, line width=1.35pt},
  inner/.style={orange!85!black, line width=1pt, line cap=round},
  outside/.style={red!65!black, line width=0.85pt},
  boundary/.style={circle, draw=blue!70!black, fill=white, inner sep=1.6pt},
  point/.style={circle, fill=blue!70!black, inner sep=1.5pt},
  guide/.style={gray!45, dashed, line width=0.4pt}
]
  \draw[axis] (-4.3,0) -- (4.4,0) node[right] {$x$};

  \draw[outside] (-4.05,0) -- (-2.25,0);
  \draw[outside] (2.25,0) -- (4.05,0);
  \draw[conv] (-2.2,0) -- (2.2,0);
  \draw[inner] (-1.35,0.48) -- (1.35,0.48);
  \draw[inner] (-1.35,0.38) -- (-1.35,0.58);
  \draw[inner] (1.35,0.38) -- (1.35,0.58);

  \node[boundary] at (-2.2,0) {};
  \node[boundary] at (2.2,0) {};
  \node[point] at (0,0) {};

  \draw[guide] (-2.2,0.08) -- (-2.2,0.58);
  \draw[guide] (2.2,0.08) -- (2.2,0.58);

  \node[below] at (-2.2,-0.05) {$a-R$};
  \node[below] at (0,-0.05) {$a$};
  \node[below] at (2.2,-0.05) {$a+R$};
  \node[above] at (0,0.62) {$[a-\rho,a+\rho]$};
  \node[blue!65!black, below] at (0,-0.42) {$\abs{x-a}<R$};
\end{tikzpicture}

\caption{Inside $\abs{x-a}<R$: absolute and locally uniform; outside divergence; endpoints undecided.}
  \label{fig:power-series-radius-interval}
\end{figure}

\begin{definition} \label{def:real-analytic-at-point} \hlt{Real Analytic at a Point}\\
Let $E\subseteq\R$, let $f:E\to\R$, and let $a$ be an interior point of $E$. The function $f$ is \hlt{real analytic at $a$} if there are $r>0$ and coefficients $(c_n)$ such that $(a-r,a+r)\subseteq E$, the power series $\sum c_n(x-a)^n$ has radius of convergence at least $r$, and $f(x)=\sum\limits_{n=0}^{\infty}c_n(x-a)^n$ for all $x\in(a-r,a+r)$.\\
If $E$ is open and $f$ is real analytic at every point of $E$, then $f$ is \hlt{real analytic on $E$}.
\end{definition}

\begin{theorem} \label{thm:power-series-uniform-differentiation} \hlt{Power Series May Be Differentiated Term by Term}\\
Let $f(x)=\sum_{n\ge0}c_nx^n$ have radius $R>0$, and $I_R=(-R,R)$. For $0<\rho<R$, let $K_\rho=[-\rho,\rho]$, a closed interval contained in $I_R$. Then $\sum c_nx^n$ converges uniformly on each $K_\rho$, and $f\in C^\infty(I_R)$.\\
For $k\ge1$, $f^{(k)}(x)=\sum\limits_{n\ge k}\frac{n!}{(n-k)!}c_nx^{n-k}$, $(x\in I_R)$, and $f^{(k)}(0)=k!c_k$.
\begin{proof}
For $0<\rho<R$, Theorem \ref{thm:power-series-radius-convergence} and the root-test argument give $\sum\abs{c_n}\rho^n<\infty$; hence the Weierstrass $M$-test, Theorem \ref{thm:weierstrass-m-test}, gives uniform convergence on $K_\rho$.\\
Since $\sqrt[n]{n}\to1$, polynomial factors in $n$ do not change the radius. Thus each differentiated series converges uniformly on $K_\rho$, and Theorem \ref{thm:uniform-derivative-convergence}, iterated, gives the formula and $f^{(k)}(0)=k!c_k$.\\
\hlt{Structural Commentary}: On each $K_\rho$, termwise operations are justified.\\
\hlt{Key Insight}: The derivative coefficients grow only by polynomial factors, which do not change the radius.
\end{proof}
\end{theorem}

\begin{corollary} \label{cor:centered-power-series-termwise-calculus} \hlt{Termwise Calculus for Centered Power Series}\\
Let $f(x)=\sum\limits_{n=0}^{\infty}c_n(x-a)^n$ have radius $R>0$. On $(a-R,a+R)$ the series converges locally uniformly and defines a continuous function. For $\abs{x-a}<R$, $f'(x)=\sum\limits_{n=1}^{\infty}n c_n(x-a)^{n-1}$ with local uniform convergence of the derivative series. If $[y,z]\subset(a-R,a+R)$, then $\int_y^z f(x)\,dx = \sum\limits_{n=0}^{\infty}c_n\frac{(z-a)^{n+1}-(y-a)^{n+1}}{n+1}$.
\begin{proof}
Set $g(u)=\sum c_nu^n$ and $u=x-a$.\\
Theorem \ref{thm:power-series-uniform-differentiation} applied to $g$ on $(-R,R)$ gives convergence, continuity, derivative formula after translating.\\
For $[y,z]\subset(a-R,a+R)$, choose $\rho<R$ with $[y,z]\subset[a-\rho,a+\rho]$; then Corollary \ref{cor:termwise-riemann-stieltjes-integration} lets the integrals of partial sums pass to the uniform limit.\\
\hlt{Structural Commentary}: All local termwise operations happen on closed intervals strictly inside the radius.\\
\hlt{Key Insight}: Translation reduces centered power series to the radius-$R$ theorem about $\sum c_nu^n$.
\end{proof}
\end{corollary}

\begin{corollary} \label{cor:analytic-functions-smooth-and-unique-taylor} \hlt{Analytic Functions Are Their Taylor Series}\\
If $f(x)=\sum c_n(x-a)^n$ near $a$, then $f$ is $C^\infty$, expansion $f$ is unique, and $f(x)=f(a)+\sum\limits_{n=1}^{\infty}\frac{f^{(n)}(a)}{n!}(x-a)^n$.
\begin{proof}
Corollary \ref{cor:centered-power-series-termwise-calculus}, applied repeatedly, gives $c_0=f(a)$ and $c_k=f^{(k)}(a)/k!$.\\
Hence two expansions about $a$ have the same coefficients.\\
\hlt{Structural Commentary}: Derivatives determine the local series.\\
\hlt{Key Insight}: Once a power-series representation exists, the coefficients cannot vary.
\end{proof}
\end{corollary}

\begin{perspective} \label{per:taylor-coefficients-need-a-remainder} \hlt{Taylor Coefficients Are Only the First Half}\\
Power-series coefficients are forced by derivatives (Corollary \ref{cor:analytic-functions-smooth-and-unique-taylor}); for a function not already known to be analytic, proving the Taylor series represents the function needs a remainder estimate (Theorem \ref{thm:taylor-theorem-lagrange-remainder}).
\end{perspective}

\begin{theorem} \label{thm:abel-endpoint-theorem} \hlt{Abel's Endpoint Theorem}\\
Let $s_n\defeq\sum_{k=0}^{n}c_k$ with $s_n\to s$, and let $f(x)=\sum_{n\ge0}c_nx^n$ on $(-1,1)$. Then $f(x)\to s$, for $x\to1^-$.
\begin{proof}
With $s_{-1}=0$, summation by parts, Theorem \ref{thm:summation-by-parts}, gives for $\abs{x}<1$, $f(x)=(1-x)\sum\limits_{n=0}^{\infty}s_nx^n$.\\
Since $(1-x)\sum\limits_{n\ge0}x^n=1$, then $f(x)-s=(1-x)\sum\limits_{n\ge0}(s_n-s)x^n$. Choose $N$ so that $\abs{s_n-s}<\varepsilon/2$ for $n>N$.\\
The tail is at most $\varepsilon/2$, and the finite part tends to $0$ after multiplication by $1-x$.\\
\hlt{Structural Commentary}: Endpoint convergence becomes one-sided continuity.\\
\hlt{Key Insight}: The factor $(1-x)$ averages the partial sums $s_n$ with geometric weights.
\end{proof}
\end{theorem}

\begin{theorem} \label{thm:abel-product-theorem} \hlt{Abel Product Theorem for Cauchy Products}\\
Suppose $\sum a_n=A$, $\sum b_n=B$, $\sum c_n=C$, and $c_n=\sum\limits_{k=0}^{n}a_kb_{n-k}$. Then $C=AB$.
\begin{proof}
For $0\leq x<1$, set $f(x)\defeq\sum\limits_{n\ge0}a_nx^n$; $g(x)\defeq\sum\limits_{n\ge0}b_nx^n$; and $h(x)\defeq\sum\limits_{n\ge0}c_nx^n$.\\
Since convergent series have bounded terms, $f,g,h$ converge absolutely for $\abs{x}<1$. Thus the Cauchy product gives $fg=h$ on $[0,1)$. By Theorem \ref{thm:abel-endpoint-theorem}, $f\to A$, $g\to B$, and $h\to C$ as $x\to1^-$. Hence $AB=C$.\\
\hlt{Structural Commentary}: Prove $fg=h$ inside, then pass to $x=1$.\\
\hlt{Key Insight}: Inside $\abs{x}<1$, the Cauchy product is safe; Abel then moves the identity to $x=1$.
\end{proof}
\end{theorem}

\begin{theorem} \label{thm:power-series-reexpansion-and-identity} \hlt{Re-Expansion and Identity for Power Series}\\
Let $f(x)=\sum_{n\ge0}c_nx^n$ have radius $R>0$, and set $S=(-R,R)$.
\begin{enumerate}[label=(\roman*)]
\setlength{\itemsep}{0pt}
\item If $a\in S$, then $f(x)=\sum\limits_{m\ge0}\frac{f^{(m)}(a)}{m!}(x-a)^m$, $(\abs{x-a}<R-\abs{a})$
\item Let $F=\sum\limits_{n\ge0}a_nx^n$ and $G=\sum\limits_{n\ge0}b_nx^n$ on $S$.\\
If $\mathcal{Z}_{F,G}\defeq\{x\in S:F(x)=G(x)\}$ has a limit point in $S$, then $a_n=b_n$ for all $n\ge0$, hence $F=G$ on $S$.
\end{enumerate}
\begin{proof}
\begin{enumerate}[label=(\roman*)]
\setlength{\itemsep}{0pt}
\item Write $x^n=((x-a)+a)^n$. If $\abs{x-a}+\abs{a}<R$, then $\sum\limits_{n\ge0}\abs{c_n}\bigl(\abs{x-a}+\abs{a}\bigr)^n<\infty$, so the resulting double sum may be rearranged. Theorem \ref{thm:power-series-uniform-differentiation} gives the coefficients $f^{(m)}(a)/m!$.
\item Subtract the series and call the difference $f$. Let $Z=\{x\in S:f(x)=0\}$, and let $A$ be the set of limit points of $Z$ in $S$. Then $A\neq\emptyset$. For $x_0\in A$, part (i) gives $f(x)=\sum d_n(x-x_0)^n$.\\
If $d_k$ were the first nonzero coefficient, then $f(x)=(x-x_0)^kg(x)$ with $g(x_0)\neq0$, contradicting $x_0\in A$. Hence $f=0$ near $x_0$, so $A$ is open. Also $S\setminus A$ is open; by connectedness, Theorem \ref{thm:connected-subsets-of-real-line-interval-property}, $A=S$. Since $Z$ is closed, $f=0$ on $S$, and evaluating derivatives at $0$ gives equality of coefficients.
\end{enumerate}
\hlt{Structural Commentary}: Power-series equality is local-to-global on an interval.\\
\hlt{Key Insight}: A limit point of zeros forces all Taylor coefficients to vanish.
\end{proof}
\end{theorem}

\begin{figure}[H]
  \centering
  \begin{tikzpicture}[
  >=Stealth,
  every node/.style={font=\scriptsize},
  x=2.55cm,
  y=0.88cm,
  axis/.style={->, black!55, line width=0.45pt},
  broad/.style={black!55, dashed, line width=0.75pt},
  narrow/.style={black!85, line width=0.9pt},
  window/.style={fill=gray!12, draw=none},
  guide/.style={black!35, dashed, line width=0.45pt}
]
  \draw[window] (-0.32,0) rectangle (0.32,2.75);

  \draw[axis] (-1.16,0) -- (1.16,0) node[right] {$t$};
  \draw[axis] (0,-0.15) -- (0,2.86);

  \draw[broad, domain=-1:1, samples=150]
    plot (\x,{1.25*(1-\x*\x)^2});
  \draw[narrow, domain=-1:1, samples=180]
    plot (\x,{2.48*(1-\x*\x)^8});

  \node at (0.66,0.88) {$Q_m$};
  \node at (0.34,1.75) {$Q_n$};

  \draw[guide] (-0.32,0) -- (-0.32,2.55);
  \draw[guide] (0.32,0) -- (0.32,2.55);
  \draw[<->, black!55] (-0.32,-0.22) -- node[below] {$2\delta$} (0.32,-0.22);

  \node[below] at (-1,0) {$-1$};
  \node[below] at (1,0) {$1$};
\end{tikzpicture}

  \caption{The normalized kernels $Q_n(t)=c_n(1-t^2)^n$ concentrate near $0$ while retaining integral $1$.}
  \label{fig:polynomial-bump-kernel}
\end{figure}

\begin{theorem} \label{thm:weierstrass-polynomial-approximation}\label{thm:power-approx-fourier-main-source-todo} \hlt{Weierstrass Polynomial Approximation}\\
If $f:[a,b]\to\C$ is continuous, then polynomials $P_n$ satisfy $P_n\rightrightarrows f$ on $[a,b]$. If $f$ is real-valued, take $P_n$ real.
\begin{proof}
By an affine change of variables and subtraction of the endpoint interpolant, assume $[a,b]=[0,1]$ and $f(0)=f(1)=0$; extend $f$ by $0$ outside $[0,1]$. For $n\geq1$, let $Q_n(x)\defeq c_n(1-x^2)^n$ and $\int_{-1}^{1}Q_n(x)\,dx=1$.\\
Source estimates give $c_n<\sqrt n$ and $Q_n(t)\leq\sqrt n(1-\delta^2)^n$ for $\delta\leq\abs{t}\leq1$.\\
Define $P_n(x)\defeq\int_{-1}^{1}f(x+t)Q_n(t)\,dt$, $(0\leq x\leq1)$. By change of variables, $P_n(x)=\int_0^1 f(u)Q_n(u-x)\,du$, so $P_n$ is a polynomial in $x$ (real if $f$ is real). Given $\varepsilon>0$, choose $\delta>0$ so $\abs{u-v}<\delta$ implies $\abs{f(u)-f(v)}<\varepsilon/2$, and set $M=\sup\abs{f}$. Since $Q_n\geq0$ and $\int Q_n=1$, $\abs{P_n(x)-f(x)} \leq \int_{-1}^{1}\abs{f(x+t)-f(x)}Q_n(t)\,dt \leq \frac{\varepsilon}{2}+4M\sqrt n(1-\delta^2)^n$.\\
The last term tends to $0$, so $P_n\rightrightarrows f$ on $[0,1]$. Undoing the reductions proves the theorem.\\
\hlt{Structural Commentary}: The kernels $Q_n$ form polynomial approximate identities.\\
\hlt{Key Insight}: Uniform continuity controls the near part, and concentration of $Q_n$ kills the far part.
\end{proof}
\end{theorem}

\begin{figure}[H]
  \centering
  \begin{tikzpicture}[
  >=Stealth,
  x=4.8cm,
  y=1.75cm,
  every node/.style={font=\scriptsize},
  axis/.style={->, gray!70, line width=0.45pt},
  curve/.style={blue!70!black, line width=0.9pt},
  coarse/.style={orange!82!black, line width=0.75pt},
  fine/.style={teal!70!black, line width=0.65pt},
  point/.style={circle, fill=black, inner sep=1.05pt},
  note/.style={fill=white, inner sep=0.8pt}
]
  \draw[axis] (-0.04,0) -- (1.06,0) node[right] {$x$};
  \draw[axis] (0,-0.06) -- (0,1.08) node[above] {$y$};

  \draw[curve, domain=0:1, variable=\x, samples=140]
    plot (\x,{sqrt(1-\x)});
  \draw[coarse] (0,1) -- (0.5,0.7071) -- (1,0);
  \draw[fine] (0,1) -- (0.125,0.9354) -- (0.25,0.8660) -- (0.375,0.7906)
    -- (0.5,0.7071) -- (0.625,0.6124) -- (0.75,0.5) -- (0.875,0.3536) -- (1,0);

  \foreach \x/\y in {0/1,0.5/0.7071,1/0} {
    \node[point] at (\x,\y) {};
  }

  \node[blue!70!black, note] at (0.30,0.93) {$f$};
  \node[orange!82!black, note] at (0.64,0.56) {$P_2$};
  \node[teal!70!black, note] at (0.82,0.42) {$P_8$};
  \node[below] at (0,0) {$0$};
  \node[below] at (1,0) {$1$};
\end{tikzpicture}

  \caption{Finer piecewise-linear interpolants approach $f(x)=\sqrt{1-x}$ uniformly on $[0,1]$.}
  \label{fig:polygonal-approximation-sqrt-one-minus-x}
\end{figure}

\begin{corollary} \label{cor:absolute-value-polynomial-approximation} \hlt{Polynomial Approximation of $\abs{x}$}\\
For every interval $[-a,a]$, there are real polynomials $P_n$ such that $P_n(0)=0$ and $P_n\rightrightarrows \abs{x}$ on $[-a,a]$.
\begin{proof}
By Theorem \ref{thm:weierstrass-polynomial-approximation}, choose real polynomials $P_n^\ast$ converging uniformly to $\abs{x}$ on $[-a,a]$.\\
Since $P_n^\ast(0)\to0$, the polynomials $P_n(x)\defeq P_n^\ast(x)-P_n^\ast(0)$ have $P_n(0)=0$ and still converge uniformly to $\abs{x}$.\\
\hlt{Structural Commentary}: Normalize the approximants at the base point.\\
\hlt{Key Insight}: Subtracting the value at $0$ fixes the normalization without changing the uniform limit.
\end{proof}
\end{corollary}

\begin{definition} \label{def:algebra-uniform-closure-functions} \hlt{Algebras and Uniform Closure}\\
Let $\mathcal{A}$ be a family of complex-valued functions on a set $E$.
\begin{enumerate}[label=(\roman*)]
\setlength{\itemsep}{0pt}
\item $\mathcal{A}$ is an \hlt{algebra} if it is closed under $+$, products, and scalar multiplication (real scalars in the real case).
\item $\mathcal{A}$ is \hlt{uniformly closed} if $f_n\in\mathcal{A}$ and $f_n\rightrightarrows f$ on $E$ imply $f\in\mathcal{A}$.
\item The \hlt{uniform closure} of $\mathcal{A}$ is $\{f:f_n\in\mathcal{A},\ f_n\rightrightarrows f\text{ on }E\}$.
\end{enumerate}
\end{definition}

\begin{theorem} \label{thm:uniform-closure-algebra} \hlt{Uniform Closures Are Closed Algebras}\\
The uniform closure of an algebra of bounded functions is a uniformly closed algebra.
\begin{proof}
If $f,g\in\mathcal{B}$, choose $f_n,g_n\in\mathcal{A}$ with $f_n\rightrightarrows f$ and $g_n\rightrightarrows g$.\\
Boundedness and uniform convergence give $f_n+g_n\rightrightarrows f+g$; $f_ng_n\rightrightarrows fg$; and $cf_n\rightrightarrows cf$.
Thus $\mathcal{B}$ is an algebra. If $h_n\in\mathcal{B}$ and $h_n\rightrightarrows h$, approximate each $h_n$ uniformly by some $a_n\in\mathcal{A}$ within $1/n$; then $a_n\rightrightarrows h$, so $h\in\mathcal{B}$.\\
\hlt{Structural Commentary}: Uniform limits preserve algebraic operations for bounded functions.\\
\hlt{Key Insight}: The uniform closure of an algebra has already absorbed all further uniform limits.
\end{proof}
\end{theorem}

\begin{definition} \label{def:separates-points-vanishes-nowhere-self-adjoint} \hlt{Separation, Nonvanishing, and Self-Adjointness}\\
Let $\mathcal{A}$ be a family of functions on $E$.
\begin{enumerate}[label=(\roman*)]
\setlength{\itemsep}{0pt}
\item $\mathcal{A}$ \hlt{separates points} on $E$ if for every distinct $x_1,x_2\in E$ there is $f\in\mathcal{A}$ with $f(x_1)\neq f(x_2)$.
\item $\mathcal{A}$ \hlt{vanishes at no point} of $E$ if for every $x\in E$ there is $g\in\mathcal{A}$ with $g(x)\neq0$.
\item A complex algebra $\mathcal{A}$ is \hlt{self-adjoint} if $f\in\mathcal{A}$ implies $\overline{f}\in\mathcal{A}$, where $\overline{f}(x)\defeq\overline{f(x)}$.
\end{enumerate}
\end{definition}

\begin{theorem} \label{thm:two-point-interpolation-in-separating-algebra} \hlt{Two-Point Interpolation in a Separating Algebra}\\
If an algebra $\mathcal{A}$ on $E$ separates points and vanishes nowhere, then for distinct $x_1,x_2\in E$ and scalars $c_1,c_2$, there is $f\in\mathcal{A}$ with $f(x_i)=c_i$, $(i=1,2)$.
\begin{proof}
Choose $g,h,k\in\mathcal{A}$ with $g(x_1)\neq g(x_2)$, $h(x_1)\neq0$, and $k(x_2)\neq0$.\\
Let $u\defeq gk-g(x_1)k$ and $v\defeq gh-g(x_2)h$. Then $u(x_1)=v(x_2)=0$ and $u(x_2)v(x_1)\neq0$.\\
Hence $f\defeq \frac{c_1v}{v(x_1)}+\frac{c_2u}{u(x_2)}$ lies in $\mathcal{A}$ and satisfies $f(x_i)=c_i$.\\
\hlt{Structural Commentary}: Separation and nonvanishing build two functions that isolate the two points.\\
\hlt{Key Insight}: Once two point-isolating functions are in the algebra, interpolation is linear algebra.
\end{proof}
\end{theorem}

\begin{perspective} \label{per:stone-weierstrass-interpolation-linear-algebra} \hlt{Interpolation Is Linear Algebra}\\
The two-point interpolation step can also be written as a $2\times2$ linear system. Its invertibility is exactly the finite-dimensional criterion from Theorem \ref{thm:finite-dimensional-operator-invertible-equivalences}, with the determinant test supplied by Theorem \ref{thm:determinant-basic-laws}. The proof above chooses explicit point-isolating functions instead of reopening the calculation.
\end{perspective}

\begin{theorem} \label{thm:real-stone-weierstrass} \hlt{Stone--Weierstrass Theorem, Real Form}\\
$\mathcal{A}\subset C(K,\R)$ is algebra on compact $K$ that separates points, vanishes nowhere, then uniform closure is $C(K,\R)$.
\begin{proof}
Let $\mathcal{B}$ be the uniform closure of $\mathcal{A}$; by Theorem \ref{thm:uniform-closure-algebra}, $\mathcal{B}$ is a uniformly closed algebra.\\
If $f\in\mathcal{B}$, then $\abs{f}\in\mathcal{B}$: for $a=\norm{f}_{\infty}$, approximate $\abs{y}$ on $[-a,a]$ by polynomials $p_n$ with $p_n(0)=0$ using Corollary \ref{cor:absolute-value-polynomial-approximation}; then $p_n(f)\in\mathcal{B}$ and $p_n(f)\rightrightarrows\abs{f}$.\\
Hence $\max(f,g),\min(f,g)\in\mathcal{B}$ whenever $f,g\in\mathcal{B}$, as $\max(f,g)=\frac{f+g}{2}+\frac{\abs{f-g}}{2}$; $\min(f,g)=\frac{f+g}{2}-\frac{\abs{f-g}}{2}$.\\
Fix $f\in C(K,\R)$, $x\in K$, $\varepsilon>0$. $\forall y\in K$, Theorem \ref{thm:two-point-interpolation-in-separating-algebra} gives $h_y\in\mathcal{B}$, $h_y(x)=f(x)$, $h_y(y)=f(y)$.\\
Compactness gives $y_1,\ldots,y_r$ s.t. $g_x\defeq\max(h_{y_1},\ldots,h_{y_r})$ satisfies $g_x(x)=f(x)$ and $g_x>f-\varepsilon$ on $K$.\\
Compactness gives $x_1,\ldots,x_m$ such that $h\defeq\min(g_{x_1},\ldots,g_{x_m})$ satisfies $\abs{h-f}<\varepsilon$ on $K$.\\
Thus every $f\in C(K,\R)$ is a uniform limit of members of $\mathcal{B}$, so $f\in\mathcal{B}$.\\
\hlt{Structural Commentary}: Local supports plus finite max/min give global approximation.\\
\hlt{Key Insight}: Separation plus compactness gives uniform approximation.
\end{proof}
\end{theorem}

\begin{theorem} \label{thm:complex-stone-weierstrass} \hlt{Stone--Weierstrass Theorem, Complex Form}\\
If $\mathcal{A}\subset C_b(K)$ is self-adjoint, separates points, and vanishes nowhere on compact $K$, then $\overline{\mathcal{A}}=C_b(K)$.
\begin{proof}
Let $\mathcal{A}_{\R}\defeq\{f\in\mathcal{A}:f\text{ real-valued}\}$.\\
If $f=u+iv\in\mathcal{A}$, self-adjointness gives $u=\frac{f+\overline{f}}{2}\in\mathcal{A}_{\R}$ and $v=\frac{f-\overline{f}}{2i}\in\mathcal{A}_{\R}$.\\
Thus $\mathcal{A}_{\R}$ separates points; rescaling a nonzero value and taking real parts shows it vanishes nowhere.\\
By Theorem \ref{thm:real-stone-weierstrass}, $\overline{\mathcal{A}_{\R}}=C(K,\R)$. Approximate the real and imaginary parts separately, then combine.\\
\hlt{Structural Commentary}: Self-adjointness lets the complex algebra expose enough real functions.\\
\hlt{Key Insight}: The complex theorem is the real theorem applied to real and imaginary parts.
\end{proof}
\end{theorem}

\begin{theorem} \label{thm:exponential-logarithm-from-power-series} \hlt{Exponential and Logarithm from the Power Series}\\
Define $\mathsf{E}(z)\defeq\sum\limits_{n=0}^{\infty}\frac{z^n}{n!}$ for $z\in\C$.\\
Then, for $z,w\in\C$, $x\in\R$, $\mathsf{E}(z+w)=\mathsf{E}(z)\mathsf{E}(w)$, $\mathsf{E}(x)>0$, $\mathsf{E}'(x)=\mathsf{E}(x)$, $\mathsf{E}(x)=e^x$, and $\mathsf{E}$ is strictly increasing on $\R$. Its inverse $\mathsf{L}:(0,\infty)\to\R$ satisfies $\mathsf{L}'(y)=\frac1y$, $\mathsf{L}(uv)=\mathsf{L}(u)+\mathsf{L}(v)$ for $u,v>0$.\\
For $x>0$ and $\alpha\in\R$, set $x^\alpha\defeq \mathsf{E}(\alpha \mathsf{L}(x))$; then $(x^\alpha)'=\alpha x^{\alpha-1}$.
\begin{proof}
The ratio test (Theorem \ref{thm:ratio-test}), gives convergence for all $z\in\C$.\\
Absolute convergence and Cauchy product give $\mathsf{E}(z)\mathsf{E}(w)=\sum\limits_{n\ge0}\frac{(z+w)^n}{n!}=\mathsf{E}(z+w)$, hence $\mathsf{E}(z)\mathsf{E}(-z)=1$.\\
Theorem \ref{thm:power-series-uniform-differentiation} gives $\mathsf{E}'=\mathsf{E}$ on $\R$. Thus $\mathsf{E}(x)>0$ $\forall$ real $x$ and $\mathsf{E}'>0$, so $\mathsf{E}$ is strictly increasing.\\
Since $\mathsf{E}(1)=e$, the addition formula gives $\mathsf{E}(p)=e^p$ for $p\in\Q$; continuity extends this to $x\in\R$.\\
Let $\mathsf{L}=\mathsf{E}^{-1}$. Differentiating $\mathsf{L}(\mathsf{E}(x))=x$ gives $\mathsf{L}'(y)=1/y$.\\
If $u=\mathsf{E}(x)$ and $v=\mathsf{E}(y)$, then $\mathsf{L}(uv)=\mathsf{L}(\mathsf{E}(x+y))=\mathsf{L}(u)+\mathsf{L}(v)$.\\
Finally, the chain rule, Theorem \ref{thm:chain-rule-one-variable}, gives $(x^\alpha)'=\alpha x^{\alpha-1}$.\\
\hlt{Structural Commentary}: Power series define $\mathsf{E}$; inversion gives $\mathsf{L}$ and real powers.\\
\hlt{Key Insight}: The addition formula is the bridge from coefficients to the logarithm law.
\end{proof}
\end{theorem}

\begin{perspective} \label{per:real-powers-rational-approximation-route} \hlt{Real Powers by Rational Approximation}\\
One route defines $x^\alpha$ for $x>0$ by choosing $q_n\in\Q$ with $q_n\to\alpha$ and setting $x^\alpha=\lim_{n\to\infty}x^{q_n}$ after proving existence and independence of $(q_n)$. The power-series construction sets $x^\alpha=\mathsf{E}(\alpha\mathsf{L}(x))$; both routes must agree with rational powers and be choice-independent.
\end{perspective}

\begin{perspective} \label{per:ross-logarithmic-integral-route} \hlt{Logarithmic Integral Route}\\
With the Riemann integral, set $L(y)\defeq\int_1^y \frac1t\,dt$ for $y>0$. Theorem \ref{thm:integral-function-continuity-and-derivative} gives $L'(y)=1/y$; endpoint limits make $L$ invertible, and Theorem \ref{thm:chain-rule-one-variable} then recovers $E'=E$ and the product laws.
\end{perspective}

\begin{perspective} \label{per:exponential-factorial-and-improper-integrals} \hlt{From Exponentials to Factorials}\\
Theorems \ref{thm:exponential-logarithm-from-power-series}, \ref{thm:riemann-integration-by-parts} supply real-power laws and integration-by-parts pattern used for $\Gamma$ in Theorem \ref{thm:gamma-function-basic-theory}.
\end{perspective}

\begin{theorem} \label{thm:trigonometric-functions-from-exponential} \hlt{Trigonometric Functions from the Exponential}\\
For real $x$, define $\mathsf{C}(x)\defeq\frac{\mathsf{E}(ix)+\mathsf{E}(-ix)}{2}$ and $\mathsf{S}(x)\defeq\frac{\mathsf{E}(ix)-\mathsf{E}(-ix)}{2i}$.\\
Then $\mathsf{E}(ix)=\mathsf{C}(x)+i\mathsf{S}(x)$, $\mathsf{C}^2+\mathsf{S}^2=1$; $\mathsf{C}'=-\mathsf{S}$, and $\mathsf{S}'=\mathsf{C}$.\\
If $x_0=\min\{x>0:\mathsf{C}(x)=0\}$ and $\pi\defeq2x_0$, then $\mathsf{E}(\pi i)=-1$, $\mathsf{E}(2\pi i)=1$, $\mathsf{E}(z+2\pi i)=\mathsf{E}(z)$ for $z\in\C$.\\
Thus $\mathsf{C},\mathsf{S}$ are $2\pi$-periodic, and $\gamma(t)\defeq \mathsf{E}(it)$ traces the unit circle once on $[0,2\pi]$.
\begin{proof}
The definitions give $\mathsf{E}(ix)=\mathsf{C}+i\mathsf{S}$. Since $\mathsf{E}(ix)\mathsf{E}(-ix)=1$, $\abs{\mathsf{E}(ix)}=1$, hence $\mathsf{C}^2+\mathsf{S}^2=1$.\\
Differentiating the definitions and using Theorem \ref{thm:power-series-uniform-differentiation} gives $\mathsf{C}'=-\mathsf{S}$ and $\mathsf{S}'=\mathsf{C}$.\\
If $\mathsf{C}$ had no positive zero, then $\mathsf{C}>0$ and $\mathsf{S}'>0$ on $(0,\infty)$.\\
Thus, for $0<x<y$, $\mathsf{S}(x)(y-x)<\int_x^y \mathsf{S}(t)\,dt=\mathsf{C}(x)-\mathsf{C}(y)\leq2$, impossible for large $y$. Hence $x_0$ exists.\\
On $(0,\pi/2)$, $\mathsf{C}>0$, so $\mathsf{S}$ increases from $0$ to $\mathsf{S}(\pi/2)=1$; thus $\mathsf{E}(\pi i/2)=i$, and the addition formula gives $\mathsf{E}(\pi i)=-1$, $\mathsf{E}(2\pi i)=1$, and the period.\\
Finally, $\abs{\mathsf{E}(it)}=1$ and the source quadrant argument give the one-time traversal of the unit circle.\\
\hlt{Structural Commentary}: Trigonometry comes from the path $\mathsf{E}(it)$.\\
\hlt{Key Insight}: The identity $\mathsf{E}(ix)\mathsf{E}(-ix)=1$ puts the exponential path on the unit circle.
\end{proof}
\end{theorem}

\begin{figure}[H]
  \centering
  \begin{tikzpicture}[
  scale=1.25,
  >=Stealth,
  every node/.style={font=\scriptsize},
  axis/.style={->, gray!70, line width=0.45pt},
  circlecurve/.style={blue!70!black, line width=0.9pt},
  arcpath/.style={orange!80!black, line width=0.9pt},
  guide/.style={gray!45, dashed, line width=0.4pt},
  point/.style={circle, fill=blue!70!black, inner sep=1.4pt},
  labelbox/.style={fill=white, fill opacity=0.9, text opacity=1, inner sep=1.2pt}
]
  \draw[axis] (-1.35,0) -- (1.5,0) node[right] {$\operatorname{Re}$};
  \draw[axis] (0,-1.35) -- (0,1.5) node[above] {$\operatorname{Im}$};

  \draw[circlecurve] (0,0) circle (1);
  \draw[arcpath, ->] (1,0) arc[start angle=0,end angle=55,radius=1];

  \coordinate (P) at (55:1);
  \draw[guide] (P) -- ({cos(55)},0) node[below] {$\mathsf{C}(t)$};
  \draw[guide] (P) -- (0,{sin(55)}) node[left] {$\mathsf{S}(t)$};
  \draw[guide] (0,0) -- (P);

  \node[point] at (P) {};
  \node[labelbox, anchor=south west] at (P) {$\mathsf{E}(it)=\mathsf{C}(t)+i\mathsf{S}(t)$};
  \node[labelbox] at (0.5,0.18) {$t$};
  \node[below right] at (1,0) {$1$};
  \node[above left] at (0,1) {$i$};
\end{tikzpicture}

  \caption{$\mathsf{E}(it)=\mathsf{C}(t)+i\mathsf{S}(t)$ traces the unit circle, defining sine and cosine analytically.}
  \label{fig:exponential-unit-circle}
\end{figure}

\begin{theorem} \label{thm:fundamental-theorem-algebra-exponential-proof} \hlt{Algebraic Completeness of $\C$}\\
Every nonconstant complex polynomial has a complex root.
\begin{proof}
Let $P(z)=\sum_{j=0}^{n}a_jz^j$, $n\geq1$, and normalize $a_n=1$.\\
Since $\abs{P(z)}\to\infty$, Theorem \ref{thm:extreme-value-theorem-compact} gives $z_0$ with $\mu\defeq\abs{P(z_0)}=\min_{\C}\abs{P}$. If $\mu=0$, done.\\
Assume $\mu>0$ and set $Q(w)\defeq P(z_0+w)/P(z_0)$. Then $Q(0)=1$ and $\abs{Q(w)}\geq1$.\\
Write $Q(w)=1+\sum\limits_{j=k}^{n}b_jw^j$, where $b_k\neq0$ and $k$ is minimal. By Theorem \ref{thm:trigonometric-functions-from-exponential}, choose $\theta$ with $b_ke^{ik\theta}=-\abs{b_k}$.\\
Then $Q(re^{i\theta})=1-\abs{b_k}r^k+r^{k+1}H(r)$, with $H$ bounded near $0$, so $\abs{Q(re^{i\theta})}<1$ for small $r>0$, contradiction.\\
\hlt{Structural Commentary}: Minimize, translate, then perturb.\\
\hlt{Key Insight}: Choose the phase that lowers the first nonzero term.
\end{proof}
\end{theorem}

\begin{theorem} \label{thm:gamma-function-basic-theory} \hlt{Gamma Function, Beta Function, and Stirling Skeleton}\\
For $x>0$, $\Gamma(x+1)=x\Gamma(x)$; $\Gamma(n+1)=n!$ for $n\in\N$; $\log\Gamma$ is convex on $(0,\infty)$.\\
Moreover, $\Gamma$ is the unique $f:(0,\infty)\to(0,\infty)$ with $f(1)=1$, $f(x+1)=xf(x)$, and $\log f$ convex.\\
For $x,y>0$, $\mathrm{B}(x,y)\defeq\int_0^1t^{x-1}(1-t)^{y-1}\,dt =\frac{\Gamma(x)\Gamma(y)}{\Gamma(x+y)}$.\\
Consequently, $\Gamma\!\left(\frac12\right)=\sqrt{\pi}$; $\int_{-\infty}^{\infty}e^{-s^2}\,ds=\sqrt{\pi}$; $\Gamma(x)=\frac{2^{x-1}}{\sqrt{\pi}}\Gamma\!\left(\frac{x}{2}\right)\Gamma\!\left(\frac{x+1}{2}\right)$; and $\lim_{x\to\infty}\frac{\Gamma(x+1)}{(x/e)^x\sqrt{2\pi x}}=1$.
\begin{proof}
Integration by parts gives $\Gamma(x+1)=x\Gamma(x)$; $\Gamma(1)=1$ gives $\Gamma(n+1)=n!$.\\
Holder's inequality gives log-convexity. For uniqueness, let $\varphi=\log f$.\\
If $0<x<1$, convexity and $\varphi(x+1)=\varphi(x)+\log x$ give $0\leq \varphi(x)-\log\!\frac{n!\,n^x}{x(x+1)\cdots(x+n)} \leq x\log\!\left(1+\frac1n\right)\to0$.\\
Thus $f$ is uniquely determined on $(0,1)$, hence on $(0,\infty)$.\\
For fixed $y>0$, $\mathrm{B}(1,y)=1/y$, $\mathrm{B}(x+1,y)=\frac{x}{x+y}\mathrm{B}(x,y)$, and Holder gives log-convexity in $x$.\\
Hence $F_y(x)\defeq\frac{\Gamma(x+y)}{\Gamma(y)}\mathrm{B}(x,y)$ satisfies the three uniqueness conditions, so $F_y=\Gamma$.\\
The consequences follow from $t=\sin^2\theta$, $t=s^2$, the duplication identity, and the substitution $t=x(1+u)$ in the asymptotic proof of Stirling's formula.\\
\hlt{Structural Commentary}: Functional equation plus convexity gives uniqueness.\\
\hlt{Key Insight}: Bohr--Mollerup uniqueness turns beta-function identities into gamma identities.
\end{proof}
\end{theorem}

\begin{method} \label{meth:power-approx-fourier-proof-patterns-source-todo}\label{meth:power-series-approximation-proof-patterns} \hlt{Power-Series and Approximation Proof Patterns}
\begin{enumerate}[label=(\roman*)]
\setlength{\itemsep}{0pt}
\item \hlt{Power series}: work on $K_\rho\Subset I_R$; use uniform convergence before termwise limits or derivatives.
\item \hlt{Endpoint theorems}: prove identities inside the radius and pass to endpoints with Abel-type control.
\item \hlt{Approximate identities}: use kernels $Q_n\geq0$, $\int Q_n=1$, with tails $\to0$ off $0$.
\item \hlt{Uniform closure}: show $a_n\in\mathcal{A}$ and $a_n\rightrightarrows f$.
\item \hlt{Lattice closure}: derive $\abs{f}$, then $\max,\min$ from algebraic formulas.
\item \hlt{Stone--Weierstrass}: interpolate locally; use compactness and finite max/min to globalize.
\item \hlt{Gamma function}: use $\Gamma(x+1)=x\Gamma(x)$, $\Gamma(1)=1$, and log-convexity for uniqueness.
\end{enumerate}
\hlt{Key Insight}: Uniform convergence and compactness are what make infinite algebra behave like finite algebra.
\end{method}
